\documentclass[11pt,a4paper]{article}
\usepackage[utf8]{inputenc}
\usepackage[T2A,T1]{fontenc}
\usepackage[russian,english]{babel}
\usepackage[margin=2.2cm]{geometry}
\usepackage{microtype}
\usepackage{xcolor}
\usepackage[colorlinks=true,urlcolor=blue]{hyperref}
\setlength{\parindent}{0pt}
\setlength{\parskip}{0.8em}
\linespread{1.05}
\pagestyle{empty}

\begin{document}

\begin{center}
{\large\bfseries Forecasting Daily Mentions of Stalin and Lenin\\[0.15em]
on Russian Television}

\vspace{0.5em}
\textbf{Team members:} Kupcha Sofya, Paltseva Ksenia
\end{center}

We plan to forecast daily mentions of Joseph Stalin and Vladimir Lenin on four Russian television channels: Channel One, NTV, Russia 1, and Russia 24. For each person and channel, we will predict two quantities: the total number of surname mentions across all programmes and the number of broadcasts containing at least one mention. Our motivation is to examine how references to these Soviet leaders reflect television rhetoric. We will produce daily point forecasts and predictive distributions for the month immediately following the last available observation.

We will use the television transcripts and broadcast metadata already collected from \href{https://gdeltproject.org/}{GDELT} through Google Cloud and stored in our project dataset. The current data cover January 2023 through October 2, 2026, so the forecast period will run from October 3 through November 2, 2026. We will count surname occurrences, including their grammatical forms, and aggregate them by broadcast, channel, and day. A broadcast containing several mentions will count once toward the number of broadcasts with mentions. We will distinguish missing transcripts from genuine zero counts and examine transcription errors that could affect the measurements.

We will compare four Bayesian count models used in our preliminary forecast: (1) a constant-rate Poisson baseline; (2) a negative-binomial model allowing greater day-to-day variation; (3) a negative-binomial regression with a time trend and partially pooled month-of-year and day-of-week effects; and (4) an extension of the third model with a latent time-varying level, estimated using a particle filter. These extensions address excess variation, calendar patterns, and changes in the underlying frequency of mentions. Each model will be evaluated at multiple historical forecast dates using an expanding-window backtest: we will fit it on earlier observations and forecast every day of the following month. We will compare daily RMSE, CRPS, and the coverage of 80\% and 95\% predictive intervals, and use Diebold--Mariano tests to assess differences in forecast performance.

\end{document}
